\section{Related Work}
\label{sec:related}

\paragraph{Cooperative Multi-Agent Reinforcement Learning.}
Decentralized cooperative MARL has predominantly developed within the Centralized Training with Decentralized Execution (CTDE) paradigm. Value-factorization methods, including Value-Decomposition Networks (VDN)~\cite{sunehag2018value} and QMIX~\cite{rashid2020monotonic}, decompose joint action-value functions under additivity or monotonicity constraints. Policy gradient frameworks have similarly advanced from Counterfactual Multi-Agent Policy Gradients (COMA)~\cite{foerster2018counterfactual}---which employs a counterfactual centralized critic to isolate agent-specific marginal contributions---to Multi-Agent PPO (MAPPO)~\cite{yu2022surprising}, which uses generalized advantage estimation with centralized value functions. Trust-region variants such as Heterogeneous-Agent Trust Region Policy Optimization (HATRPO) and HAPPO~\cite{kuba2022trust} formalize theoretical multi-agent monotonic improvement bounds. In role-oriented MARL, architectures such as ROMA~\cite{wang2020roma} and RODE~\cite{wang2021rode} construct emergent dynamic role embeddings or decompose action spaces via clustering. However, existing CTDE and role-based methods maintain static network capacities; when deployed in environments with temporally mismatched reward frequencies, shared parameters succumb to gradient interference, where dense early rewards corrupt sparse late-stage representations.

\paragraph{Mixture of Experts and Dynamic Capacity Growth.}
Sparsely-gated Mixture-of-Experts (MoE) architectures~\cite{shazeer2017outrageously,fedus2022switch,lepikhin2020gshard} scale model capacity without proportional inference computation by dispatching inputs to subsets of specialized sub-networks via learned gating. In continual learning and transfer, structural expansion approaches dynamically grow network capacity when loss stagnates: Progressive Neural Networks~\cite{rusu2016progressive} instantiate lateral feature columns, Dynamically Expandable Networks (DEN)~\cite{yoon2018lifelong} dynamically add neurons upon high task loss, and DEMO~\cite{yan2021dynamically} expands MoE pools across lifelong task sequences. However, existing dynamic capacity techniques trigger expansion exclusively upon explicit, external task boundaries or rigid heuristic schedules. In contrast, THIEF introduces \emph{deficit-driven} expert pool expansion operating entirely within a single long-horizon task, using empirical performance plateaus and failure transitions as the internal expansion signal.

\paragraph{Gradient Conflict and Multi-Task Optimization.}
When an agent or shared network simultaneously optimizes disparate objectives, parameter updates often encounter severe multi-task gradient conflict. Yu et al.~\cite{yu2020gradient} formalized this pathology through Projected Conflicting Gradients (PCGrad), demonstrating that negative cosine similarity between task gradient vectors ($\cos(\nabla \mathcal{L}_i, \nabla \mathcal{L}_j) < 0$) causes destructive cross-task interference. While PCGrad projects gradients onto normal planes to mitigate interference in fixed architectures, THIEF repurposes the cosine-similarity conflict diagnostic between successful and failed trajectory partitions as an autonomous triggering signal for architectural mitosis: rather than forcing a single network to compromise across conflicting behavioral manifolds, THIEF spawns a dedicated specialist to absorb the conflicting sub-manifold.

\paragraph{Model Merging and Fisher Information Geometry.}
Combining diverse neural models without accessing full historical training datasets is a foundational problem in model merging. Kirkpatrick et al.~\cite{kirkpatrick2017overcoming} utilized the diagonal of the empirical Fisher Information Matrix (FIM) in Elastic Weight Consolidation (EWC) to constrain updates away from parameter directions critical to historical tasks. Matena and Raffel~\cite{matena2021merging} mathematically formalized Fisher-weighted averaging, proving that merging fine-tuned models weighted by their diagonal Fisher information yields curvature-optimal parameter combinations on the underlying statistical manifold. While Fisher-merging has been investigated as a post-hoc operation for combining independently trained large models, THIEF leverages this Riemannian geometry as an \emph{online genetic recombination operator}, synthesizing parent capabilities into newly spawned specialists during active RL training.

\paragraph{Evolutionary Reinforcement Learning and Population-Based Training.}
Evolutionary Strategies (ES)~\cite{salimans2017evolution,stanley2019open} and Evolutionary Reinforcement Learning (ERL)~\cite{khadka2018evolution} combine gradient-free parameter exploration with policy-gradient sample efficiency. Population-Based Training (PBT)~\cite{jaderberg2017population} continuously adapts hyperparameters and policy weights through periodic exploit-and-explore cycles across parallel model instances. THIEF adapts the spirit of PBT's population dynamics but departs from random parameter mutation: THIEF's sandbox warmup, cooldown gating, and dormancy pruning establish a disciplined lifecycle where specialist generation is strictly targeted at empirical failure deficits, calibrated in isolated environment shards, and governed by switching barriers.
